\xchapter{Hymni Ecclesiæ Pars I—E Breviario Parisiensi}

\xsection{Preface}

OF the three kinds of poetical composition which, in accordance
with the Apostle's direction, have ever been in use in the Church, ``Psalms
and Hymns and Spiritual Songs,'' two are supplied by inspiration. We
have no need, through God's bounty, to turn our thoughts to the
composition of Psalms or Songs; and, to judge from the attempts which
have been made, doubtless we are unequal to it. And the unapproachable
excellence of the two which have been supplied serves to suggest the
difficulties which beset the composition of the third which has not
been supplied. Indeed, it is hardly too strong to say that to write
Hymns is as much beyond us as to originate Psalmody. The peculiarity
of the Psalms is their coming nearer than any other kind of devotion
to a converse with the powers of the unseen world. They are longer and
freer than Prayers; and, as being so, are less a direct address
to the Throne of Grace than a sort of intercourse, first with oneself,
then with one's brethren, then with Saints and Angels, nay, even the
world and all creatures. They consist mainly of the praises of God;
and the very nature of praise involves a certain abstinence from
intimate approaches to Him, and the introduction of other beings into
our thoughts, through whom our offering may come round to Him. For as
He, and He only, is the direct object of prayer, so it is more
becoming not to regard Him as directly addressed in praise, which
would imply passing a judgment on Him who is above all scrutiny and
all standards. The Seraphim cried \emph{one to another}, ``Holy, Holy,
Holy,'' veiling their faces; neither looking nor speaking to Him. The
Psalms, then, as being praises, and thanksgivings, are the language,
the ordinary converse, as is may be called, of Saints and Angels in
heaven; and, being such, could not be written except by men who had
heard the ``unspeakable things'' which there are uttered. In this light
they are more difficult than Prayers. Beggars can express their wants
to a prince; they cannot converse like his courtiers.

Much the same remark may be made about the Songs
or Canticles of the Church, which are also inspired, and are a kind of
Psalms written for particular occasions, chiefly occasions of
thanksgiving. Such are the two Songs of Moses, the Song of Hannah,
those in Isaiah, the Song of Hezekiah, of Habakkuk, of the Three
Children, of Zacharias, of the Blessed Virgin, and of Simeon; most of
which are in the Breviary, and the last four are retained in our own
Reformed Prayer Book.

Yet though Hymns, as being of a measured length,
and restrained metrically, are so far safer to attempt than Psalms or
Songs, they have their own peculiar difficulties. They are direct
addresses to Almighty God, which ever must be most difficult to the
serious mind, whatever be the difficulty of other devotions. This, in
the instance of Prayers, has led to the use of Sentences, such as
occur in our own Services; which, besides the advantage of extreme
brevity, for the most part admit of being taken from Scripture. It has
led also to the repetition of the Lord's Prayer, and of the \emph{Kyrie
Eleison}; and, again, to the use of Collects, which lessen the
difficulty of addressing God by subjecting it to fixed rules. Hence
our best Family Prayers are what may be called a succession of
Sentences strung together, the simple and concise expression of our
humiliation, fear, hope, and desire for ourselves and others. Long
Public Prayers, to make a general assertion which of course admits of
exceptions, are arrogant and irreverent; hence the Pharisees made
them. Hence, too, the unchastised effusions which abound in the
present day among those who have left the Church or lost her spirit.
The great Eucharistic Prayer is nearly the only long prayer in the Catholic Church; and there is every reason to suppose that in its
substance it proceeds from inspired authorities. In our own Service it
has been separated by our Reformers into three distinct portions.

Hymns, however, being of the nature of praises,
cannot be altogether brought down to that grave and severe character
which, as being direct addresses to God, they seem to require; and
this is their peculiar difficulty. To praise God specially for
Redemption, to contemplate the mysteries of the Divine Nature, to
enlarge upon the details of the Economy of Grace and yet not to
offend, to invoke with awe, to express affection with a pure heart, to
be subdued and sober while we rejoice, and to make professions without
display, and all this not under the veil of figurative language, as in
the Psalms, but plainly, and (as it were) abruptly, surely requires to
have had one's lips touched with a ``coal from the Altar,'' to have
caught from heaven that ``new song'' ``which no man could learn, but the
hundred and forty and four thousand which were redeemed from the earth''—the
virgin followers of the Lamb.

Our Church, with the remarkable caution which she
displays so often, has not attempted it. She has received the Psalms
and Songs from Scripture; and, rejecting the Roman Hymns, has
substituted in their stead, not others, but a metrical version of the
Psalms. This abstinence has led on one hand to some of her
members on their own responsibility supplying the deficiency, and has
incurred the complaint of others, who argued that she ought to have
taken on herself what, being right in itself, will certainly be done
by private hands, if not by the fitting authority. But, in truth, when
it was necessary for her to abandon those she had received, nothing
was left to her but to wait till she should receive others, as in the
course of ages she had already received, by little and little.

The Roman Hymns, whether good or bad, were the
work of no one generation, much less the outpourings of one mind. They
were not the contents of one collection, published all new in a day
according to the will of man. They were the gradual accumulations of
centuries, bearing in old and new upon one treasure-house. When there
was a call to reject them, there was nothing to be done but begin
again. We could not be young and old at once. It was a stern necessity
alone which could compel us to change from what we were; but bang
changed, so far we were not what we were, and must be what the
primitive Church was in these respects, poor and ill-furnished. We
began the world again. This is the proper answer to inconsiderate
complaints and impatient interference. There have before now been
divines who could write a Liturgy in thirty-six hours. Such is not our
Church's way. She is not the empiric to make things to order, and to
profess and to anticipate the course of nature, which, under grace, as
under Providence, is slow. She waits for that majestic course to
perfect in its own good time, what she cannot extort from it; for the
gradual drifting of precious things upon her shore, now one and now
another, out of which she may complete her rosary and enrich her
beads,—beads and rosary more pure and true than those which at the
command of duty she flung away.

As far as we know, the public Hymns of the early
Church were not much more than the following. First, starting from
Scripture, she adopted the repetition of the \emph{Hallelujah}, which
is described by St. John, in the Revelations, to be the chant of the
blessed inhabitants of heaven. Next may be mentioned the \emph{Gloria
Patri}, pretty much as we now use it. Thirdly, the \emph{Trisagion},
or ``Holy, Holy, Holy,'' from Isaiah vi.; or, as it was also used, and
now is, in the Roman Church, ``Sanctus Deus, Sanctus Fortis, Sanctus
Immortalis.'' Besides these, there was the Morning or Angelic Hymn,
beginning with the words used by the Angels at the nativity; and for
the evening the Hymn beginning ``Hail gladdening Light,'' preserved by
St. Basil. These are not metrical, as they were afterwards; nor are
two others of a later date, which we still retain, the \emph{Te Deum}
and the \emph{Athanasian Creed}. They are both of Gallican origin,
though the former has been ascribed to St. Ambrose. Others, however,
now extant, are certainly his; others are the compositions of St.
Hilary, Prudentius, St. Gregory, and later saints. It is not
too much to say then that, judging by what we know of the Hymns of the
primitive Church, we should not be dissatisfied with the paucity of
those which custom has, with a sort of tacit authority, introduced
among us in the course of several centuries.

More, doubtless, might be selected from the
writings of our sacred Poets; but since, from unhappy circumstances,
such a work does not seem likely at the present day, thoughtful minds
naturally revert to the discarded collections of the ante-reform era,
discarded because of associations with which they were then viewed,
and of the interpolations by which they were disfigured; but which,
when purified from these, are far more profitable to the Christian
than the light and wanton effusions which are their present substitute
among us. Nay, even such as the Parisian, which are here first
presented to the reader, which have no equal claims to antiquity,
breathe an ancient spirit; and even where they are the work of one
pen, are the joint and invisible contribution of many ancient minds.
Moreover, the ancient language used has a tendency to throw the reader
out of every-day thoughts and familiar associations, and to make him
fervent without ceasing to be mortified. Many a man could bear to read
the Canticles in a foreign language who is unequal to it in his own.

It only remains to say, that the following
selection of Hymns, from the Paris Breviary, has been confined to such
holy days and seasons as are recognized by our Church, or to
special events or things recorded in Scripture; those Hymns, however,
being omitted which contained invocations to the Saints of such a
nature as to be, even in the largest judgment of charity, not mere
apostrophes, but supplications.
J. H. N.

February 21, 1838.

[Thanks to an anonymous reader for calling this preface to our attention, and to Brian
Cardell, Assistant Curator, Rare Books and Special Collections,
Catholic University of America Libraries, for providing copies of
parts I and II from the 1865 edition—\textbf{NR}.]

\xsection{Pars II—E Breviariis Romano, Sarisburiensi, Eboracensi et aliunde}

\xsection{Præfatio}

NE mirum tibi, Lector benevole, videatur, Hymnos, quos in manibus
habes, et materie vetustos et stilo, hodiernis typis et hodiernorum
favori essi commendatos. Non enim abs te exigitur, ut quicquid hic
scriptum inveneris, id in sinum tuum recipias illico, absque animi
ulla exercitatione aut judicio tuo. Neque qui totum aliquid edit in
publicum, ilium oportet universi operis in singula quæque jurare, aut
in rerum squalore immorari quod in antiquitate versatur. Alia nimirum
sunt in monumentis veterum simplicis ac pristinæ naturæ elementa;
prorsus aliud concretum vitium, et sordes temporum. Cordati autem
hominis est, necue obsoleta pro priscis, nec peregrina pro falsis
habere.

Quod quidem optimo consilio actum est olim a
patribus nostris, qui, cæco quodam reformationis, quam vocant,
æstu in Ecclesia passim fervente, e religionibus suis, non quæ
exterorum studium posceret, sed quæ vis extorqueret, missa fecerunt;
a malesana philosophia ereptis cæteris, ea concedendo, quæ
sanctitatem plane nullam, parem autem invidiam haberent. Hinc illi,
quibus nunc utimur, ritus et cæremoniæ; hinc rerum divinarum ordo;
hinc et discipinæ nostræ forma et agendi ratio repetenda est.

Quod si fiat quæstio, cur diligentissimi viri,
cum nostra omnia retractarent, hymnodiam solum intactam, et tanquam
desperatam, ne dicam improbatam reliquerint, id vero ex ejusmodi
silentio minime colligendum erit. Nam, cum statutum esset, ut cultus
divinus, eo usque in Latino idiomate absolutus, transiret in
vernaculam, quis daret, ut qui possent aliquid in rebus reformandis,
ii in poetica florerent simul, verba autem in versum et quæstiones ad
veritatis normam, pari solertia redigerent? Unde factum est, ut in
officiis divinis, dum habemus cætera, hymnos non habeamus.

Quare in hac Ecclesiæ nostræ qualicunque inopia,
hoc tempore, Lector benevole, tibi apponenda statuimus illa vel
illorum similia, quæ nostri in sæculo decimo sexto, cum sua essent,
de manibus prudenter amiserunt; non illa, ex omni parte
approbanda, sed aspersa maculis, ita vero ut modo vix quicquam
improbes, modo improbes, at neque ut rejicias nec tamen ut possis
mederi. Nam quæ tota reprobanda errant, ista scilicet tota omisimus;
cætera subjicimus judicio tuo.

J. H. N.

\xsection{Title Page}

HYMNI ECCLESIAE

\emph{PARS I}

E
BREVIARIO PARISIENSI

\emph{PARS II}

E BREVIARIIS

ROMANO, SARISBURIENSI, EBORACENSI

ET ALIUNDE

Londini

APUD

ALEXANDRUM MACMILLAN

1865
